\section{Introdución}
El presente informe, documenta el desarrollo del proyecto \enquote{Control de Acceso para institución educativa} desarrollado en el marco de la asignatura de Proyecto II a modo de proyecto final de carrera de Ingeniería Electrónica.
El proyecto consta de un sistema de control de acceso de dos vías, una convencional para personal autorizado mediante identificadores personales de readiofrecuencia, y otra para visitantes externos los cuales al solicitar un acceso tocando el timpre de la puerta, son fotografiados para que la imagen sea enviada a un teléfono celular donde la persona a cargo validará el acceso y permitirá o no la apertura de la puerta.
El diseño se utiliza tomando como base la Escuela de Enseñanza Secundaria Obligatoria N° 214 "Mariano Moreno", ubicada en la localidad de Santa Isabel, la cual cuenta con tres puntos de acceso que requieren de este sistema.
Sobre cada puerta se colocará un llamado "Terminal" basado en la placa de desarrollo ESP32-CAM y se comunicarán con un nodo central llamdo "Servidor" ubicado en una oficina. La interfaz del usuario que valida los accesos mediante un dispositivo movil se ejecutará mediante un bot de la aplicación de mensajería instantanea "Télegram".

